
Ovarian cancer is the leading cause of gynecological cancer death and is the seventh most common cancer in women worldwide.
Epithelial ovarian cancer comprises about 90\% of all ovarian malignancies and includes high grade serous, low grade serous, mucinous, endometrioid and clear cell ovarian carcinomas.
High grade serous ovarian carcinomas (HGSOC) represent more than half of all epithelial ovarian cancers and is the most widely studied epithelial ovarian tumor, while endometrioid and mucinous carcinomas account for 10\% and 6\%, respectively \cite{torre2018}.
Through extensive whole genome, exome, and methylation analyses of HGSOC, a variety of genes have been identified as likely drivers in these tumors \cite{cancergenomeatlasresearchnetwork2011}, and both clinical and evolutionary analyses have implicated the fallopian tube as the tissue of origin of HGSOC \cite{labidi-galy2017, eckert2016, ducie2017, piek2001, finch2006, medeiros2006, kuhn2012, perets2013, rabban2014}.
%Alice - adding piek2001, finch2006, medeiros2006, kuhn2012, perets2013, and rabban2014 as clinico-pathologic studies of fallopian tube HGSOC origin.
Apart from alterations in a few well-known drivers including \textit{KRAS} and \textit{HER2} amplification in ovarian serous and mucinous carcinomas and mutations in \textit{ARID1A}, \textit{PIK3CA}, and \textit{PTEN} genes as well as mismatch repair (MMR) deficiency in ovarian endometrioid carcinomas, comparatively little is known about molecular alterations in these tumors \cite{slamon1989}.
Comprehensive genomic analyses of these tumor subtypes could delineate additional alterations and pathways frequently altered in these cancers and shed light on their likely tissue of origin.
% Nothing about uterine / GI mucinous cancers in first paragraph.

Uterine endometrioid adenocarcinomas account for nearly 85\% of all epithelial uterine carcinomas and accumulating evidence suggests a shared molecular etiology between these cancers and ovarian endometrioid adenocarcinomas \cite{buhtoiarova2016}.
Synchronous endometrioid ovarian and uterine carcinomas can arise from independent primary tumors or as metastases from a single primary tumor \cite{makris2017, hajkova2019}, suggesting that these tumors may share the same cell of origin \cite{wu2001, wu2017}.
Women with synchronous ovarian and uterine cancers have improved survival compared to women with a single cancer, in part due to early diagnosis from the symptoms of vaginal bleeding associated with uterine carcinomas \cite{matsuo2018}. %is this relevant to our goal?
Understanding the origin of these cancers could lead to more accurate prognoses of survival or more aggressive early intervention.
While whole genome, whole exome, and targeted next generation sequencing analyses have been performed on ovarian endometrioid carcinomas \cite{wang2017, hoang2015, hollis2020},  and whole exome studies have been performed on uterine endometroid carcinomas \cite{cancergenomeatlasresearchnetwork2013}, the extent to which these cancers share similar alterations and a potentially common tissue of origin has not been systematically evaluated.

Distinguishing between primary mucinous ovarian cancer and mucinous adenocarcinoma metastases from other sites, including GI mucinous tumors, may be challenging.
Some histological features can favor primary diagnosis but there are several discordant and overlapping features that make a definitive diagnosis difficult \cite{brown2014}.
Molecular analyses for mucinous ovarian and GI cancers, including pancreatic, stomach and colorectal cancer, have largely been limited to targeted sequencing \cite{lee2017, khan2018, nakata2002, cheasley2019, dundr2023}, gene expression \cite{elias2018}, or immunohistochemistry \cite{elias2014}, though recently work has begun on exomic genomic analyses of these malignancies \cite{cheasley2019, mueller2018}.
No studies to date have compared the epigenomes of mucinous carcinomas from the ovary with those of other sites.
Such analyses could establish the tissue of origin and lead to better risk stratification and treatment for patients with these cancers.
In this study, we characterize the molecular landscapes of endometrioid and mucinous adenocarcinomas, comparing genomes and epigenomes in both the cancer and adjacent normal tissues (Figure \ref{fig:fig1}a).
Our analyses indicate the major pathways affected by these cancers and identify molecular alterations that may improve our understanding of their tissue of origin.